%%% FIELD proof-uniform-margin-band
Fix a sample outcome in \(\mathcal B_n\), and write \(r_0=r_{H,n}\).  By Definition \(\mathrm{def:simultaneous\text{-}band}\), for every \(a\in\mathcal A\) and \(u\in\Omega\),
\[
 \lVert\widetilde H_a(u)-H_a(u)\rVert_{\mathrm{op}}\le r_0.
 \tag{1}
\]
Assumption \(\mathrm{ass:hessian\text{-}ellipticity}\) gives \(H_a(u)\succeq mI_d\).  Thus, for every unit vector \(x\in\mathbb R^d\), (1) and the assumed inequality \(r_0\le m/2\) give
\[
 x^{\mathsf T}\widetilde H_a(u)x
 \ge x^{\mathsf T}H_a(u)x-r_0
 \ge m-r_0
 \ge \frac m2>\frac m4.
 \tag{2}
\]
The field \(\widetilde H_a\) is symmetric by Definition \(\mathrm{def:regularized\text{-}hessian\text{-}estimator}\).  Consequently every eigenvalue in (2) exceeds the clipping threshold \(m/4\), so the definition of \(\Pi_{m/4}\) yields
\[
 \widehat H_a(u)=\widetilde H_a(u),
 \qquad
 \lVert\widehat H_a-H_a\rVert_\infty\le r_0,
 \qquad
 \widehat H_a(u)\succeq\frac m2I_d.
 \tag{3}
\]
In particular, the spectral theorem, (3), and \(H_a(u)\succeq mI_d\) imply
\[
 \lVert\widehat H_a(u)^{-1}\rVert_{\mathrm{op}}\le\frac2m,
 \qquad
 \lVert H_a(u)^{-1}\rVert_{\mathrm{op}}\le\frac1m.
 \tag{4}
\]
This proves all three simultaneous Hessian assertions in the theorem.

Fix an ordered pair \(p\ne q\).  Adding and subtracting \(H_q\widehat H_p^{-1}\), and using the exact identity
\[
 \widehat H_p^{-1}-H_p^{-1}
 =\widehat H_p^{-1}(H_p-\widehat H_p)H_p^{-1},
 \tag{5}
\]
we obtain, pointwise in \(u\),
\[
\begin{aligned}
 \lVert\widehat H_q\widehat H_p^{-1}-H_qH_p^{-1}\rVert_{\mathrm{op}}
 &\le
 \lVert\widehat H_q-H_q\rVert_{\mathrm{op}}
 \lVert\widehat H_p^{-1}\rVert_{\mathrm{op}}\\
 &\quad+
 \lVert H_q\rVert_{\mathrm{op}}
 \lVert\widehat H_p^{-1}\rVert_{\mathrm{op}}
 \lVert H_p-\widehat H_p\rVert_{\mathrm{op}}
 \lVert H_p^{-1}\rVert_{\mathrm{op}}\\
 &\le
 \left(\frac2m+\frac{2M}{m^2}\right)r_0
 =w_n.
\end{aligned}
 \tag{6}
\]
Here submultiplicativity was used in the first inequality; (3)--(4) and the upper ellipticity bound \(\lVert H_q(u)\rVert_{\mathrm{op}}\le M\) from \(\mathrm{ass:hessian\text{-}ellipticity}\) were used in the second; and the last equality is Definition \(\mathrm{def:simultaneous\text{-}band}\).  For every square matrix \(A\), the triangle inequality and transpose invariance of the operator norm give
\[
 \lVert\operatorname{sym}(A)\rVert_{\mathrm{op}}
 \le\frac12\{\lVert A\rVert_{\mathrm{op}}+\lVert A^{\mathsf T}\rVert_{\mathrm{op}}\}
 =\lVert A\rVert_{\mathrm{op}}.
 \tag{7}
\]
Definitions \(\mathrm{def:cross\text{-}hessian\text{-}margin}\) and \(\mathrm{def:regularized\text{-}hessian\text{-}estimator}\), followed by (6)--(7), therefore give
\[
 \sup_{u\in\Omega}
 \lVert\widehat S_{pq}(u)-S_{pq}(u)\rVert_{\mathrm{op}}
 \le w_n.
 \tag{8}
\]

We next derive the needed eigenvalue inequality.  If \(A\) and \(B\) are symmetric and \(x_B\) is a unit eigenvector associated with \(\lambda_{\min}(B)\), the Rayleigh representation gives
\[
 \lambda_{\min}(A)
 \le x_B^{\mathsf T}Ax_B
 \le\lambda_{\min}(B)+\lVert A-B\rVert_{\mathrm{op}}.
\]
Interchanging \(A\) and \(B\) proves
\[
 \lvert\lambda_{\min}(A)-\lambda_{\min}(B)\rvert
 \le\lVert A-B\rVert_{\mathrm{op}}.
 \tag{9}
\]
Moreover, for any real functions \(a,b\) on the same nonempty set,
\[
 \left\lvert\inf_u a(u)-\inf_u b(u)\right\rvert
 \le\sup_u\lvert a(u)-b(u)\rvert,
 \tag{10}
\]
because \(a(u)\ge b(u)-\sup_v\lvert a(v)-b(v)\rvert\) for every \(u\), and conversely after interchanging \(a,b\).  Apply (9) pointwise in (8), then apply (10), and use Definition \(\mathrm{def:cross\text{-}hessian\text{-}margin}\) and the definition of \(\widehat\gamma_{pq}\).  The result is
\[
 \lvert\widehat\gamma_{pq}-\gamma_{pq}\rvert\le w_n.
 \tag{11}
\]
The definition \(I_{pq,n}=[\widehat\gamma_{pq}-w_n,\widehat\gamma_{pq}+w_n]\) turns (11) into \(\gamma_{pq}\in I_{pq,n}\).

Finally, if \(\gamma_{pq}\le0\), then (11) implies \(\widehat\gamma_{pq}\le w_n\); the strict-threshold definition \(\Psi_n=\mathbf 1\{\widehat\gamma_{pq}>w_n\}\) therefore gives \(\Psi_n=0\).  If \(\gamma_{pq}\ge3w_n\), then (11) gives \(\widehat\gamma_{pq}\ge2w_n>w_n\), where \(w_n>0\) follows from \(m,M,r_{H,n}>0\).  Hence \(\Psi_n=1\).  Every step was pointwise on \(\mathcal B_n\), and no probability assertion about that event was used.

%%% FIELD statement-conditional-uniform-margin-band-narrowed
Fix \(\delta\in(0,1/4)\).  Suppose the two conclusions posed in \(\mathrm{oeq:global\text{-}endpoint\text{-}response}\) are resolved affirmatively with constants uniform over \(\mathfrak E_s\): for \(0<\varepsilon\le1/2\),
\[
 \lVert D^2\varphi_g-D^2\varphi_f\rVert_\infty
 \le C\varepsilon\{1+\log(C/\varepsilon)\},
 \qquad
 \lVert\varphi_g-\varphi_f-\mathscr L_fe\rVert_{C^2(\Omega)}
 \le C\lVert e\rVert_{C^r(D_p)}^2,
 \tag{A}
\]
with the first inequality interpreted as zero when \(\varepsilon=0\).  Suppose also that the following boundary-adapted implementation estimates hold uniformly over \(\mathbf P\in\mathcal M_s\), all regimes, and all sufficiently large \(n_\ast\), with one finite auxiliary constant \(C_{\mathrm{aux}}>0\) depending only on \(d,s,R,m,M,b,B,L\), and suppose the constant in Definition \(\mathrm{def:two\text{-}resolution\text{-}envelope}\) satisfies
\[
 C\ge5\max\{C_{\mathrm{aux}},C_{\mathrm{aux}}^2\}.
 \tag{C}
\]

First, a mass-preserving biorthogonal boundary-wavelet projector gives an event \(\mathcal G_n\) such that
\[
 \inf_{\mathbf P\in\mathcal M_s}\mathbb P_{\mathbf P}(\mathcal G_n)
 \ge1-\frac\delta2,
 \tag{W1}
\]
and, on \(\mathcal G_n\), simultaneously for every \(p\in\mathcal A\),
\[
 \widehat f_p=\widetilde f_p,
 \qquad
 \lVert\widehat f_p-f_p\rVert_{C^r(D_p)}\le C_{\mathrm{aux}}\eta_n,
 \qquad
 \lVert Q_{k_n}f_p-f_p\rVert_\infty\le C_{\mathrm{aux}}k_n^s,
 \qquad
 \lVert Q_{k_n}f_p-f_p\rVert_{C^r(D_p)}\le C_{\mathrm{aux}}k_n^{s-r}.
 \tag{W2}
\]
Second, with \(B_{h_n}v=D^2(K_{h_n}v)\), a bounded extension and polynomial reproduction give, for every \(v\in C^2(\Omega)\cap H^2(\Omega)\),
\[
 \lVert B_{h_n}\varphi_p-H_p\rVert_\infty\le C_{\mathrm{aux}}h_n^s,
 \quad
 \lVert B_{h_n}v\rVert_\infty\le C_{\mathrm{aux}}\lVert v\rVert_{C^2(\Omega)},
 \quad
 \lVert B_{h_n}v\rVert_\infty\le C_{\mathrm{aux}}h_n^{-d/2}\lVert v\rVert_{H^2(\Omega)},
 \quad
 \lVert\nabla B_{h_n}v\rVert_\infty\le C_{\mathrm{aux}}h_n^{-d/2-1}\lVert v\rVert_{H^2(\Omega)}.
 \tag{K}
\]
Third, on \(\mathcal G_n\), if \(e_p=\widehat f_p-f_p\), the normalized response satisfies
\[
 \widehat\varphi_p-\varphi_p
 =\mathscr L_{f_p}e_p+\operatorname{Rem}_{f_p}(e_p),
 \quad
 \lVert\operatorname{Rem}_{f_p}(e_p)\rVert_{C^2(\Omega)}\le C_{\mathrm{aux}}\eta_n^2,
 \quad
 \lVert B_{h_n}\mathscr L_{f_p}(Q_{k_n}f_p-f_p)\rVert_\infty
 \le C_{\mathrm{aux}}k_n^s\{1+\log(1/k_n)\}.
 \tag{R}
\]
Fourth, define the centered influence kernel
\[
 G_{p,h_n,k_n}(u,y)
 =B_{h_n}\mathscr L_{f_p}\!\left[Q_{k_n}(\delta_y-P_p)\right](u).
\]
For all matrix-coordinate indices \(\iota,\jmath\in\{1,\ldots,d\}\), its entries satisfy
\[
 \sup_{u\in\Omega}\mathbb E_{P_p}
   \lvert G_{p,h_n,k_n}(u,Y_{p,i})_{\iota\jmath}\rvert^2\le C_{\mathrm{aux}}h_n^{-d},
 \quad
 \sup_{u,y}\lVert G_{p,h_n,k_n}(u,y)\rVert_{\mathrm{op}}
 \le C_{\mathrm{aux}}(h_nk_n)^{-d/2},
 \tag{I1}
\]
and
\[
 \sup_{u\ne v,\,y}
 \frac{\lVert G_{p,h_n,k_n}(u,y)-G_{p,h_n,k_n}(v,y)\rVert_{\mathrm{op}}}
      {\lVert u-v\rVert_2}
 \le C_{\mathrm{aux}}h_n^{-1}(h_nk_n)^{-d/2}.
 \tag{I2}
\]
The resulting matrix-Bernstein and source-net event \(\mathcal Z_n\) obeys
\[
 \inf_{\mathbf P\in\mathcal M_s}\mathbb P_{\mathbf P}(\mathcal Z_n)
 \ge1-\frac\delta2,
 \tag{I3}
\]
and, on \(\mathcal Z_n\), simultaneously in \(p\),
\[
 \sup_{u\in\Omega}\left\lVert
 \frac1{n_p}\sum_{i=1}^{n_p}G_{p,h_n,k_n}(u,Y_{p,i})
 \right\rVert_{\mathrm{op}}
 \le C_{\mathrm{aux}}\left\{
 \sqrt{\frac{\ell_n}{n_\ast h_n^d}}
 +\frac{\ell_n}{n_\ast(h_nk_n)^{d/2}}
 \right\}.
 \tag{I4}
\]
Then the raw Hessian envelope is class-uniform:
\[
 \inf_{\mathbf P\in\mathcal M_s}
 \mathbb P_{\mathbf P}\!\left\{
 \max_{p\in\mathcal A}\lVert\widetilde H_p-H_p\rVert_\infty
 \le\mathfrak R_{H,n}\right\}\ge1-\delta.
 \tag{B1}
\]
At \(h_n=(\ell_n/n_\ast)^{1/(2s+d)}\) and \(k_n=h_n\ell_n^{-2/s}\), there are \(C_1<\infty\) and \(n_0<\infty\), with \(n_0\) allowed to depend on the fixed \(\delta\), such that, for \(n_\ast\ge n_0\),
\[
 \mathfrak R_{H,n}\le r_{H,n}=C_1h_n^s\le\frac m2,
 \qquad
 \inf_{\mathbf P\in\mathcal M_s}\mathbb P_{\mathbf P}(\mathcal B_n)\ge1-\delta.
 \tag{B2}
\]
Consequently, for every ordered pair \(p\ne q\),
\[
 \inf_{\mathbf P\in\mathcal M_s}
 \mathbb P_{\mathbf P}\{\gamma_{pq}\in I_{pq,n}\}\ge1-\delta,
\]
and the declared-pair test has size at most \(\delta\) over \(\gamma_{pq}\le0\) and power at least \(1-\delta\) over \(\gamma_{pq}\ge3w_n\).  The selected upper radius has order
\[
 w_n\asymp
 \left\{\frac{\log(C_0Rn_\ast/\delta)}{n_\ast}\right\}^{s/(2s+d)}.
\]
The premises (A), (W1)--(W2), (K), (R), and (I1)--(I4) are a conditional analytic and stochastic program, not membership assumptions on \(\mathcal M_s\).

%%% FIELD proof-conditional-uniform-margin-band
Assume all premises in the proposed statement.  On \(\mathcal G_n\), (W2) makes the fallback inactive.  Since \(\widetilde f_p=Q_{k_n}P_{p,n_p}\) by Definition \(\mathrm{def:regularized\text{-}hessian\text{-}estimator}\), linearity of \(Q_{k_n}\) gives
\[
 e_p=\widehat f_p-f_p
 =Q_{k_n}(P_{p,n_p}-P_p)+(Q_{k_n}f_p-f_p).
 \tag{1}
\]
Apply \(B_{h_n}\) to the response expansion (R), use linearity of \(\mathscr L_{f_p}\), and recall that \(\widetilde H_p=B_{h_n}\widehat\varphi_p\).  We obtain the exact decomposition
\[
\begin{aligned}
 \widetilde H_p-H_p
 &=\{B_{h_n}\varphi_p-H_p\}
 +B_{h_n}\mathscr L_{f_p}(Q_{k_n}f_p-f_p)\\
 &\quad+B_{h_n}\mathscr L_{f_p}Q_{k_n}(P_{p,n_p}-P_p)
 +B_{h_n}\operatorname{Rem}_{f_p}(e_p).
\end{aligned}
 \tag{2}
\]
The empirical-measure definition and the definition of the centered influence kernel identify the third term as
\[
 B_{h_n}\mathscr L_{f_p}Q_{k_n}(P_{p,n_p}-P_p)
 =\frac1{n_p}\sum_{i=1}^{n_p}G_{p,h_n,k_n}(\,\cdot\,,Y_{p,i}).
 \tag{3}
\]
The first inequality in (K), the last inequality in (R), and (I4) bound the first three terms of (2).  The \(C^2\)-operator inequality in (K) and the remainder inequality in (R) give the separate, composed estimate
\[
 \lVert B_{h_n}\operatorname{Rem}_{f_p}(e_p)\rVert_\infty
 \le C_{\mathrm{aux}}^2\eta_n^2.
 \tag{4}
\]
Set \(C_\dagger=\max\{C_{\mathrm{aux}},C_{\mathrm{aux}}^2\}\).  On \(\mathcal G_n\cap\mathcal Z_n\), equations (2)--(4) imply, simultaneously over \(p\in\mathcal A\),
\[
\begin{aligned}
 \lVert\widetilde H_p-H_p\rVert_\infty
 \le C_\dagger\Bigg\{&h_n^s+k_n^s[1+\log(1/k_n)]
 +\sqrt{\frac{\ell_n}{n_\ast h_n^d}}\\
 &+\frac{\ell_n}{n_\ast(h_nk_n)^{d/2}}+\eta_n^2\Bigg\}
 \le\mathfrak R_{H,n}.
\end{aligned}
 \tag{5}
\]
The last inequality uses (C), hence in particular \(C\ge C_\dagger\), and Definition \(\mathrm{def:two\text{-}resolution\text{-}envelope}\).  Notice that (4) is exactly why the strengthened condition (C) is needed.

For each \(\mathbf P\in\mathcal M_s\), the union bound, (W1), and (I3) give
\[
 \mathbb P_{\mathbf P}(\mathcal G_n\cap\mathcal Z_n)
 \ge1-\mathbb P_{\mathbf P}(\mathcal G_n^{\mathsf c})
       -\mathbb P_{\mathbf P}(\mathcal Z_n^{\mathsf c})
 \ge1-\delta.
 \tag{6}
\]
Combining (5)--(6), then taking the infimum over \(\mathbf P\in\mathcal M_s\), proves (B1).

We next simplify the deterministic envelope at the selected resolutions.  Put \(A_\delta=C_0R/\delta\).  The restrictions \(C_0>1\), \(R\ge2\), and \(0<\delta<1/4\) imply \(A_\delta>1\), and
\[
 \ell_n=\log(A_\delta n_\ast)\asymp\log n_\ast,
 \qquad
 h_n\longrightarrow0.
 \tag{7}
\]
The definition of \(h_n\) is equivalent to
\[
 \frac{\ell_n}{n_\ast}=h_n^{2s+d},
 \qquad
 \sqrt{\frac{\ell_n}{n_\ast h_n^d}}=h_n^s.
 \tag{8}
\]
Since \(k_n=h_n\ell_n^{-2/s}\), equations (7)--(8) yield
\[
 \log(1/k_n)=\log(1/h_n)+\frac2s\log\ell_n=O(\ell_n)
 \tag{9}
\]
and therefore
\[
 k_n^s\{1+\log(1/k_n)\}
 =h_n^s\ell_n^{-2}\{1+\log(1/k_n)\}
 =O(h_n^s/\ell_n),
 \tag{10}
\]
while
\[
 \frac{\ell_n}{n_\ast(h_nk_n)^{d/2}}
 =h_n^{2s}\ell_n^{d/s}=o(h_n^s).
 \tag{11}
\]
For (11), after division by \(h_n^s\) the remaining factor is \(h_n^s\ell_n^{d/s}\), which tends to zero by (7).

Substitution of \(k_n=h_n\ell_n^{-2/s}\) and \(\ell_n/n_\ast=h_n^{2s+d}\) into the definition of \(\eta_n\) gives
\[
 \eta_n
 =h_n^{s-r}\ell_n^{-2(s-r)/s}
 +h_n^{s-r}\ell_n^{(d+2r)/s}
 +h_n^{2s-r}\ell_n^{2(d+r)/s}.
 \tag{12}
\]
Assumptions \(\mathrm{ass:dimension}\) and \(\mathrm{ass:smoothness\text{-}index}\) state \(d\ge2\) and \(s>d/2+3\), so \(s>4\).  From the declared definition \(r=(s+4)/4\),
\[
 s-2r=\frac{s-4}{2}>0.
 \tag{13}
\]
Using \((x+y+z)^2\le3(x^2+y^2+z^2)\), divide the square of (12) by \(h_n^s\).  The resulting three upper-bound terms are
\[
 h_n^{s-2r}\ell_n^{-4(s-r)/s},
 \qquad
 h_n^{s-2r}\ell_n^{2(d+2r)/s},
 \qquad
 h_n^{3s-2r}\ell_n^{4(d+r)/s}.
 \tag{14}
\]
The powers of \(h_n\) in (14) are positive by (13).  More explicitly, for every fixed \(\beta>0\) and finite \(K\), (7) gives
\[
 h_n^\beta\ell_n^K
 =n_\ast^{-\beta/(2s+d)}
   \ell_n^{K+\beta/(2s+d)}\longrightarrow0.
 \tag{15}
\]
Equations (14)--(15) prove
\[
 \eta_n^2=o(h_n^s).
 \tag{16}
\]
Equations (8), (10), (11), and (16), inserted into Definition \(\mathrm{def:two\text{-}resolution\text{-}envelope}\), show that there is a finite constant \(C'\), depending only on the fixed displayed constants and the fixed \(\delta\), such that
\[
 \mathfrak R_{H,n}\le C'h_n^s
 \tag{17}
\]
for all sufficiently large \(n_\ast\).  Choose the fixed radius constant \(C_1\ge C'\).  Since \(m>0\) and \(h_n^s\to0\), there is a finite \(n_0\), allowed to depend on the fixed \(\delta\), such that for \(n_\ast\ge n_0\),
\[
 \mathfrak R_{H,n}\le r_{H,n}=C_1h_n^s\le\frac m2.
 \tag{18}
\]
By Definition \(\mathrm{def:simultaneous\text{-}band}\), the event in (B1) is contained in \(\mathcal B_n\) when (18) holds.  Thus (B1) proves the probability assertion in (B2).

On \(\mathcal B_n\), Theorem \(\mathrm{thm:uniform\text{-}margin\text{-}band}\) applies by (18).  Its deterministic conclusions give \(\gamma_{pq}\in I_{pq,n}\) for every ordered pair.  For the declared pair, they also show that rejection under \(\gamma_{pq}\le0\), or failure to reject under \(\gamma_{pq}\ge3w_n\), can occur only on \(\mathcal B_n^{\mathsf c}\).  The probability assertion in (B2) therefore proves the stated coverage, size, and power bounds.  Finally, Definition \(\mathrm{def:simultaneous\text{-}band}\) and the definitions of \(h_n\) and \(\ell_n\) give the exact equality
\[
 w_n=\left(\frac2m+\frac{2M}{m^2}\right)C_1
 \left\{\frac{\log(C_0Rn_\ast/\delta)}{n_\ast}\right\}^{s/(2s+d)}.
\]
Its prefactor is finite and strictly positive, proving the asserted order.

%%% FIELD proof-conditional-matching-separation
Assume the complete conditional analytic and stochastic program stated in \(\mathrm{prop:conditional\text{-}uniform\text{-}margin\text{-}band}\), including its corrected envelope-constant condition, and specialize that proposition to \(\delta=0.05\).  Since \(0.05\in(0,1/4)\), the proposition supplies a finite threshold \(n_0\) such that, whenever \(n_\ast\ge n_0\), the declared-pair test \(\Psi_n\) obeys
\[
 \sup_{\mathbf P\in\mathcal M_s:\,\gamma_{pq}\le0}
 \mathbb P_{\mathbf P}(\Psi_n=1)\le0.05
 \tag{1}
\]
and
\[
 \sup_{\mathbf P\in\mathcal M_s:\,\gamma_{pq}\ge3w_n}
 \mathbb P_{\mathbf P}(\Psi_n=0)\le0.05<0.20.
 \tag{2}
\]
Thus the measurable test \(\psi=\Psi_n\) satisfies both error constraints in Definition \(\mathrm{def:testing\text{-}frontier}\) at the positive separation \(\rho=3w_n\).  Taking the infimum over the feasible positive separations in that definition gives
\[
 \rho_n^\ast\le3w_n.
 \tag{3}
\]

The definitions of \(w_n\), \(r_{H,n}\), \(h_n\), and \(\ell_n\), at \(\delta=0.05\), give
\[
\begin{aligned}
 3w_n
 &=3\left(\frac2m+\frac{2M}{m^2}\right)r_{H,n}\\
 &=3\left(\frac2m+\frac{2M}{m^2}\right)C_1h_n^s\\
 &=3\left(\frac2m+\frac{2M}{m^2}\right)C_1
 \left\{\frac{\log(C_0Rn_\ast/0.05)}{n_\ast}\right\}^{s/(2s+d)}.
\end{aligned}
 \tag{4}
\]
Every division in (4) is valid because \(m>0\), \(n_\ast\ge1\), and \(0.05>0\); the exponent is positive because \(s>0\) and \(d\ge2\).  Set \(A=C_0R/0.05\).  Since \(C_0>1\) and \(R\ge2\), one has \(A>1\).  For \(n_\ast\ge\max\{2,A\}\),
\[
 \log(C_0Rn_\ast/0.05)=\log A+\log n_\ast\le2\log n_\ast.
 \tag{5}
\]
Combining (3)--(5) yields
\[
 \rho_n^\ast
 \le
 3\left(\frac2m+\frac{2M}{m^2}\right)C_1
 2^{s/(2s+d)}
 \left(\frac{\log n_\ast}{n_\ast}\right)^{s/(2s+d)}.
 \tag{6}
\]
The prefactor is finite and independent of \(n_\ast\); denoting it, or any larger fixed value, by \(C\) proves the asserted upper inequality for all sufficiently large \(n_\ast\).

Now additionally suppose that the existing class constants admit the strict-slack clause in Definition \(\mathrm{def:fixed\text{-}support\text{-}lower\text{-}family}\).  The already established Theorem \(\mathrm{thm:sharp\text{-}separation\text{-}lower}\) supplies \(c>0\) and a finite threshold \(n_1\) such that
\[
 \rho_n^\ast\ge c
 \left(\frac{\log n_\ast}{n_\ast}\right)^{s/(2s+d)}
 \qquad(n_\ast\ge n_1).
 \tag{7}
\]
For \(n_\ast\ge\max\{n_0,n_1,2,A\}\), equations (6)--(7) prove the displayed two-sided frontier.  The upper argument used the complete conditional program in \(\mathrm{prop:conditional\text{-}uniform\text{-}margin\text{-}band}\).  The lower argument invoked only \(\mathrm{thm:sharp\text{-}separation\text{-}lower}\), so it is unconditional relative to the endpoint-response, wavelet, kernel, influence, increment, and concentration program.
